% Section 10 -- Conclusion
\section{Conclusion}
\label{sec:conclusion}

We set out to build a system that answers two questions about a BGP
anomaly---did it affect traffic, and which way did the causation
run---and to do so in a way an operator could audit afterwards. The first
question has a categorical answer. The second has an answer only when the
evidence identifies an order, and our contribution is a directional test
that says so explicitly: calibrated size, measured power, and a named
reason every time it declines to commit.

\textbf{What the system does.} BGP Tracer collects RIPE RIS updates and
Cloudflare Radar traffic into a shared five-minute
grid~\cite{ripe-ris,cloudflare-radar}, runs four detectors that emit
typed evidence with provenance, and arbitrates them into a five-class
correlation verdict under six precedence rules (\S\ref{sec:method},
\S\ref{sec:impl}). On the curated evaluation set the Phase-1
arbitration reports 76.5\,\% Top-1. The directional layer is reported
through the quantities that can be re-run end to end---its size, its power
against planted mechanisms, and its coverage on the caches in the tree
(\S\ref{sec:eval-direction})---and it commits on none of the cached
windows that yield two usable series, because every one of them stops at
the first veto and names the missing edge. Abstention therefore buys an
auditable record of what the evidence lacked, which is the property the
design targets.

\textbf{What the coordination layer does.} A shared Evidence Blackboard
with write-once semantics carries typed records---with their sources,
confidences, and the identifiers they were derived from---between
specialised agents, and the layer is decomposed so that a component
failure degrades the artefact without altering the verdict: across four
injected failures the verdict was unchanged and the real measurements
survived (\S\ref{sec:agent}). Against a monolithic single-prompt arm on
identical evidence over a 17-event ground-truth corpus, the layer buys
both accuracy---$0.76$ against $0.06$---and auditability: it answers
$5.5$ of six operator questions from typed records where the monolithic arm
answers $1.4$ (\S\ref{sec:agent-explain}). Against a pure threshold rule
the accuracy gain is absent ($0.65$), so the layer's measurable
contribution lies in arbitration over conflicting detectors and in the
artefact it produces, not in the number of its stages
(Appendix~\ref{app:agent-ablation}). One agent closes a genuine
model-driven decision--action--observation loop; the rest are fixed
pipeline stages, and we classify them by control flow rather than by
assertion.

\textbf{What the case study shows.} The Orange Spain window contains a
detector firing without its baseline, a traffic verdict computed against a
baseline that does not exist, and a synthesized record standing in for an
observation (\S\ref{sec:case}). Each would have survived a review that
read only verdicts. Together with the input-shape defect it exposes
(\S\ref{sec:eval-shape}), they argue for three-state recording and for
treating abstention as a first-class output.

\textbf{Future work.} Traffic-side evidence is the binding constraint: a
second vantage point, or an operator's own traffic split, would test
whether the $R\to T$ gap is a property of the data or of our
single-series measurement. Leak attribution needs cross-domain
confirmation, because the aggregate traffic series does not move when
paths shift between peers. And the abstention record is itself a dataset:
every refusal carries the veto that failed and the edge that was
missing---a supervision signal for a channel that could learn when it is
about to be wrong.
